\chapter{Bordism Functoriality and Extended Topological Field Theory}
\label{ch:tqft-bordism-functoriality-extended}

\ConventionsEuclidean

The cohomological mechanism of the preceding chapter produces
metric-independent correlation functions under explicit Ward-identity
hypotheses.  A functorial topological field theory packages a stronger set of
gluing laws.  This chapter states the target structure before examples such
as BF theory and Chern--Simons theory are introduced.

\section{Bordism Categories}

Fix a dimension \(D\) and a tangential structure \(\xi\), such as orientation,
spin, or framing.  The bordism category \(\operatorname{Bord}_D^\xi\) has
closed \((D-1)\)-manifolds with \(\xi\)-structure as objects and compact
\(D\)-dimensional bordisms with compatible \(\xi\)-structure as morphisms.
Disjoint union makes it a symmetric monoidal category:
\[
  \Sigma_1\sqcup\Sigma_2,\qquad
  M_1\sqcup M_2.
\]
Composition is gluing along the identified boundary component.

\begin{definition}[Atiyah--Segal topological field theory]
\label{def:atiyah-segal-tqft}
A \(D\)-dimensional topological field theory with tangential structure
\(\xi\) and linear target is a symmetric monoidal functor
\[
  Z:\operatorname{Bord}_D^\xi\longrightarrow {\bf Vect}_{\Bbbk}
\]
to vector spaces over a field \(\Bbbk\).  Thus
\[
  Z(\Sigma_1\sqcup\Sigma_2)\simeq Z(\Sigma_1)\otimes Z(\Sigma_2),
  \qquad
  Z(M_2\circ M_1)=Z(M_2)Z(M_1).
\]
\end{definition}

The vector space \(Z(\Sigma)\) is the state space assigned to the spatial
slice \(\Sigma\).  For a closed \(D\)-manifold \(M\), the value \(Z(M)\) is a
linear map \(\Bbbk\to\Bbbk\), hence a number.  The cylinder
\(\Sigma\times[0,1]\) maps to the identity on \(Z(\Sigma)\).

This is the metric-independent specialization of the functorial viewpoint
developed in Chapter~\ref{ch:kontsevich-segal-functorial-qft}.  A
Kontsevich--Segal theory keeps admissible complex metrics and analytic
dependence as part of the bordism datum; an Atiyah--Segal TQFT is obtained
when the dependence on that metric datum is constant after the required
topological anomaly data have been included.

\section{Pairing and Reflection}

If \(\bar\Sigma\) denotes \(\Sigma\) with reversed tangential structure, a
bordism from \(\varnothing\) to \(\bar\Sigma\sqcup\Sigma\) gives a bilinear
pairing
\[
  Z(\bar\Sigma)\otimes Z(\Sigma)\to\Bbbk .
\]
In unitary theories the target is a category of Hilbert spaces or a
reflection-positive variant of it.  The pairing is then compatible with the
adjoint operation obtained by reflecting bordisms.  Positivity is additional
structure on the functor, parallel to reflection positivity in Euclidean QFT.

\section{Extended Theories}

An extended TQFT assigns data to manifolds of all codimensions.  The target is
then a symmetric monoidal higher category \(\mathcal C\), and the theory is a
symmetric monoidal higher functor
\[
  Z:\operatorname{Bord}_{D}^{\xi,{\rm ext}}\to\mathcal C .
\]
For example, in a fully extended two-dimensional theory, a point is assigned
an object \(A\), an interval a morphism built from \(A\), and a surface a
linear map obtained by composing those lower-codimension data.  Locality is
encoded by the ability to cut and reglue at every codimension.

\begin{quotedtheorem}[Cobordism hypothesis, framed fully extended form]
\label{thm:cobordism-hypothesis-boundary}
Let \(\mathcal C\) be a symmetric monoidal \((\infty,D)\)-category with the
dual and adjoint data needed to evaluate framed bordisms.  An object
\(X\in\mathcal C\) is fully \(D\)-dualizable when \(X\) has a dual, the
evaluation and coevaluation \(1\)-morphisms admit left and right adjoints,
those adjunction data admit the required adjoints one categorical level
higher, and this iterated adjointability continues through level \(D\).
For \(D=1\) this is ordinary dualizability; for \(D=2\) it is
dualizability plus adjoints for the evaluation and coevaluation
\(1\)-morphisms.  Under these hypotheses, fully extended framed
\(D\)-dimensional topological field theories with target \(\mathcal C\) are
classified by fully \(D\)-dualizable objects of \(\mathcal C\), with
equivalence understood as equivalence of symmetric monoidal higher functors.
\end{quotedtheorem}

This is the framed fully extended cobordism hypothesis in Lurie's
\((\infty,n)\)-categorical formulation.  It is used here as a boundary
statement, not as a result proved in this monograph.  Its role in the present
chapter is to fix the kind of target data that an extended TQFT claim must
supply.  Whenever an example is claimed to define an extended theory, the
target category, dualizability, anomaly, and boundary data must be supplied.
The quoted source lineage is the Baez--Dolan cobordism-hypothesis conjecture,
Lurie's proof and formulation for fully extended framed theories, and modern
low-dimensional and derived-geometric expositions such as the work of
Calaque--Scheimbauer.  The proof is not reproduced here; the statement is
kept to fix the exact data an extended TQFT claim must provide.

\section{From Local QFT to Functorial Data}

A local QFT with a gapped phase can produce a topological theory of its
infrared protected sector.  The passage has several ingredients:
identification of the protected observables; proof that correlation functions
depend only on topology and background \(\xi\)-structure; finite-dimensional
state spaces after imposing the topological sector; and gluing compatibility
for cutting manifolds.  In gauge-theoretic examples this also includes gauge
fixing, BV integration cycles, anomaly cancellation, and boundary conditions.

\begin{openproblem}[Functorial extraction from QFT]
\label{op:functorial-extraction-from-qft}
Given a local QFT with a proposed topological sector, construct the associated
bordism functor directly from its Hilbert spaces, local algebras, or
Euclidean functional integrals.  The construction should identify the precise
completion of state spaces, the gluing topology, anomaly line, and boundary
category.
\end{openproblem}

\section{Four-Dimensional Cohomological Gauge Theory As A Test Case}

Witten--Donaldson theory will be a principal test case for the relation
between cohomological path integrals, functorial TQFT data, and ordinary
four-dimensional gauge theory.  Its ultraviolet presentation is the
topological twist of \(\mathcal N=2\) Yang--Mills theory.  The Donaldson side
is organized by the moduli space of anti-self-dual connections, its
compactification and orientation data, and the observables obtained from the
descent of gauge-invariant polynomials.  These are not merely formal symbols:
to define the invariant one must state the gauge group and global form, the
four-manifold and its orientation data, the metric dependence and chamber
structure, the compactification of the instanton moduli space, and the
treatment of reducible connections.

The Seiberg--Witten explanation compares this ultraviolet cohomological gauge
theory to an Abelian monopole effective theory on the Coulomb branch.  The
physical mechanism is an RG statement: after flowing from the non-Abelian
\(\mathcal N=2\) theory to its low-energy Wilsonian effective action, the
singular loci where monopoles or dyons become light supply the Abelian fields
and monopole equations that enter the low-energy topological path integral.
For the monograph this comparison is not certified by naming the flow.  A
complete argument must construct the Wilsonian coordinate \(u\), the effective
coupling \(\tau(u)\), the duality frame near each singular locus, the
monopole or dyon hypermultiplet, the contact terms and gravitational
couplings, the observable map from Donaldson classes to low-energy
observables, and the wall-crossing or \(u\)-plane terms when they are present.

\begin{openproblem}[Donaldson--SW RG comparison]
\label{op:donaldson-seiberg-witten-rg-comparison}
Develop the Donaldson/Seiberg--Witten relation as a theorem-level QFT
comparison.  The target is to identify which parts follow from rigorous
differential geometry of instanton and monopole moduli spaces, which parts
follow from finite-dimensional localization or gluing, and which parts depend
on a still-unproved RG theorem connecting the twisted non-Abelian
\(\mathcal N=2\) Yang--Mills theory to the Abelian monopole effective theory.
The desired output is an object-level map between correlation functions or
cohomological observables, with stated metric, chamber, compactification,
contact-term, and anomaly data.
\end{openproblem}
